\documentclass[11pt]{article}

\usepackage[margin=1in]{geometry}
\usepackage{amsmath, amssymb}
\usepackage{xcolor}
\usepackage{listings}
\usepackage{hyperref}
\usepackage{array}
\usepackage{booktabs}

\hypersetup{
    colorlinks=true,
    linkcolor=blue!60!black,
    urlcolor=blue!60!black,
    citecolor=blue!60!black
}

\lstdefinestyle{bugpython}{
  language=Python,
  basicstyle=\ttfamily\small,
  keywordstyle=\color{blue!60!black},
  stringstyle=\color{green!40!black},
  commentstyle=\color{gray!70!black},
  showstringspaces=false,
  breaklines=true,
  frame=single,
  rulecolor=\color{gray!30},
  columns=fullflexible,
  keepspaces=true
}

\title{SymPy Bug Report}
\author{}
\date{}

\begin{document}

\maketitle

\section{Bug 34: Branch Cut Error in Radical Solving}

\subsection{Minimal Reproducer}

\medskip

\begin{lstlisting}[style=bugpython]
from sympy import Eq, S, sqrt, solveset, symbols

x = symbols("x")
sol = solveset(Eq(sqrt(1/x), 1/sqrt(x)), x, domain=S.Complexes)

print("solution:", sol)
print("contains -1:", sol.contains(S.NegativeOne))
print("lhs at -1:", sqrt(1/x).subs(x, -1))
print("rhs at -1:", (1/sqrt(x)).subs(x, -1))
print("residual at -1:", (sqrt(1/x) - 1/sqrt(x)).subs(x, -1))
\end{lstlisting}

\subsection{Observed Behavior}

\medskip

\begin{lstlisting}[style=bugpython]
solution: Complement(Complexes, {0})
contains -1: True
lhs at -1: I
rhs at -1: -I
residual at -1: 2*I
\end{lstlisting}

SymPy reports all nonzero complex numbers as solutions, so the returned set includes \(x=-1\). Direct substitution into the original equation gives a nonzero residual, \(2i\), which proves that the reported set contains invalid solutions.

\subsection{Expected Behavior}

The returned complex-domain solution set should exclude \(0\) and the negative real axis. In particular, \(\operatorname{contains}(-1)\) should be false. A mathematically correct description is
\[
    \mathbb{C} \setminus (-\infty,0],
\]
where \(0\) is excluded because the original expression is undefined and the negative real axis is excluded by the principal-square-root branch.

\subsection{Mathematical Explanation}

SymPy's \(\sqrt{\cdot}\) denotes the principal square root. For \(x=-1\),
\[
    \sqrt{1/x} = \sqrt{-1} = i,
    \qquad
    \frac{1}{\sqrt{x}} = \frac{1}{\sqrt{-1}} = \frac{1}{i} = -i.
\]
Thus
\[
    \sqrt{1/x} - \frac{1}{\sqrt{x}} = 2i
    \quad\text{at}\quad x=-1,
\]
so \(x=-1\) is not a solution.

More generally, write a nonzero complex number as \(x = r e^{i\theta}\) with the principal argument \(\theta \in (-\pi,\pi]\). For \(\theta \in (-\pi,\pi)\), the reciprocal has principal argument \(-\theta\), and the two sides agree. On the negative real axis, however, \(\theta=\pi\) and \(1/x\) is also on the negative real axis with principal argument \(\pi\). The principal square root therefore picks the opposite sign from \(1/\sqrt{x}\). This is not a harmless alternate representation: SymPy's own substitution shows distinct values \(i\) and \(-i\).

\subsection{Source-Level Diagnosis}

The diagnosis identifies the relevant source location as \nolinkurl{sympy/solvers/solveset.py}, in the helper \texttt{\_solve\_radical} around lines 1090--1139. The traced call path is
\[
    \texttt{solveset}
    \;\to\;
    \texttt{\_solveset}
    \;\to\;
    \texttt{unrad}
    \;\to\;
    \texttt{\_solve\_radical}.
\]
For this equation, \texttt{solveset} first rewrites the input as \(\sqrt{1/x}-1/\sqrt{x}\). After denominator clearing, \texttt{unrad} reduces the numerator \(\sqrt{x}\sqrt{1/x}-1\) to an identity, recorded as \texttt{(0, [])}. The radical-free problem is then treated as true for all complex numbers except denominator singularities, producing \(\mathbb{C}\setminus\{0\}\).

The source-level issue is that \texttt{\_solve\_radical} validates finite candidate sets with \texttt{checksol}, but its infinite-set handling does not recheck the original branch-sensitive equation over the broad result. The diagnosis highlights this pattern:

\begin{lstlisting}[style=bugpython]
elif isinstance(solutions, Complement):
    A, B = solutions.args
    return Complement(check_set(A), B)
...
else:
    return solutions
\end{lstlisting}

In this case, \texttt{check\_set(S.Complexes)} remains \texttt{S.Complexes}, and only \(0\) is subtracted as a denominator singularity. The negative real axis survives even though the original principal-square-root equation rejects it. A plausible fix direction supported by the diagnosis is to preserve an explicit condition or solve the missing branch constraints whenever radical removal yields broad infinite sets such as \(\mathbb{C}\), intervals, or complements.

\subsection{Additional Failing Instances}

The independent verification checks the representative value \(x=-1\) and five additional negative real values. For each case, SymPy's returned set claims membership while direct principal-square-root evaluation gives a nonzero residual. The expected behavior is the same across the family because every negative real value lies on the principal square-root branch cut, where \(\sqrt{1/x}\) and \(1/\sqrt{x}\) choose opposite signs.

\begin{center}
\small
\renewcommand{\arraystretch}{1.3}
\begin{tabular}{@{}>{\raggedright\arraybackslash}p{0.44\linewidth}>{\raggedright\arraybackslash}p{0.23\linewidth}>{\raggedright\arraybackslash}p{0.23\linewidth}@{}}
\toprule
\textbf{Input / Expression} & \textbf{SymPy behavior} & \textbf{Expected behavior} \\
\midrule
\(x=-1\) in \(\sqrt{1/x}=1/\sqrt{x}\) & Returned set contains \(-1\); residual \(2i\) & Exclude \(-1\) \\
\(x=-2\) in \(\sqrt{1/x}=1/\sqrt{x}\) & Returned set contains \(-2\); residual nonzero & Exclude \(-2\) \\
\(x=-3\) in \(\sqrt{1/x}=1/\sqrt{x}\) & Returned set contains \(-3\); residual nonzero & Exclude \(-3\) \\
\(x=-1/2\) in \(\sqrt{1/x}=1/\sqrt{x}\) & Returned set contains \(-1/2\); residual nonzero & Exclude \(-1/2\) \\
\(x=-5\) in \(\sqrt{1/x}=1/\sqrt{x}\) & Returned set contains \(-5\); residual nonzero & Exclude \(-5\) \\
\(x=-7/3\) in \(\sqrt{1/x}=1/\sqrt{x}\) & Returned set contains \(-7/3\); residual nonzero & Exclude \(-7/3\) \\
\bottomrule
\end{tabular}
\end{center}

\subsection{Related Issues Assessment}

The bug-finding harness performed a best-effort deduplication check and classified this as a new instance of a known failure mode. The closest related public material concerns the general principal-square-root branch-cut family: SymPy documentation explains that \(\sqrt{\cdot}\) is principal, SymPy simplification guidance warns that power identities need assumptions, and public mathematical discussions cover why square-root multiplicativity or reciprocal identities fail on branch cuts. The available evidence did not identify a clear public duplicate for this exact \texttt{solveset} call returning \(\mathbb{C}\setminus\{0\}\) and including \(x=-1\). The case is therefore individually actionable as a concrete solver regression test and fix target.

\subsection{Proposed SymPy Regression Test}

\medskip

\begin{lstlisting}[style=bugpython]
import pytest

from sympy import Eq, Rational, S, simplify, sqrt, solveset, symbols


@pytest.mark.parametrize(
    "bad_value",
    [
        S.NegativeOne,
        S(-2),
        S(-3),
        Rational(-1, 2),
        S(-5),
        Rational(-7, 3),
    ],
)
def test_solveset_sqrt_reciprocal_excludes_negative_real_axis(bad_value):
    x = symbols("x")
    lhs = sqrt(1/x)
    rhs = 1/sqrt(x)

    sol = solveset(Eq(lhs, rhs), x, domain=S.Complexes)

    assert sol.contains(bad_value) is S.false
    assert simplify((lhs - rhs).subs(x, bad_value)) != 0
\end{lstlisting}

\end{document}
